We saw the \textbf{XOR problem} as an example of how linear models limits, 

\section{Notation}
To maintain consistency throughout this document, we adopt the standard 
notation summarized below, largely following standard deep learning 
literature:
\begin{table}[h]
\centering
\small
\begin{tabular}{ll}
\toprule
\textbf{Symbol} & \textbf{Description} \\
\midrule
$x \in \mathbb{R}^d$ & one input, with $d$ features \\
$y$ (or $\mathbf{y}$) & its target: a number, a class, or a vector \\
$f^*$ & the unknown true function, $y = f^*(x)$ \\
$n$ & number of training examples \\
$(x^{(i)}, y^{(i)})$ & the $i$-th training example \\
$X \in \mathbb{R}^{n \times d}$ & all inputs, one example per row \\
$m$ & minibatch size (examples per update) \\
$k$ & number of classes \\
\bottomrule
\end{tabular}
\caption{Data notation summary.}
\label{tab:data_notation}
\end{table}

\begin{table}[h]
\centering
\small
\begin{tabular}{ll}
\toprule
\textbf{Symbol} & \textbf{Description} \\
\midrule
$L$ & number of layers (the depth) \\
$d_l$ & number of units in layer $l$ (its width); $d_0 = d$ \\
$W^{(l)}, b^{(l)}$ & weight matrix ($d_{l-1} \times d_l$) and bias vector of layer $l$ \\
$z^{(l)}$ & pre-activation: $W^{(l)\top} h^{(l-1)} + b^{(l)}$ \\
$h^{(l)} = g(z^{(l)})$ & activation of layer $l$; $h^{(0)} = x$ \\
$g$ & activation function of hidden units \\
$o$ & output function (identity, sigmoid or softmax) \\
$\theta$ & all weights and biases together \\
\bottomrule
\end{tabular}
\caption{Network notation summary.}
\label{tab:network_notation}
\end{table}

\noindent\textit{Note on indexing:} Superscripts in parentheses, such as 
$^{(l)}$ or $^{(i)}$, always denote indices (referring to layers or training
examples, respectively) and are never used as mathematical powers. 
Vectors are represented in bold lowercase letters, matrices in bold 
uppercase, and scalars in standard font.

\begin{table}[htbp]
\centering
\small
\begin{tabular}{ll}
\toprule
\textbf{Symbol} & \textbf{Description} \\
\midrule
$\hat{y} = f(x; \theta)$ & prediction of the network \\
$\ell(\hat{y}, y)$ & loss on one example \\
$J(\theta)$ & cost: average loss on the training set \\
$\nabla_\theta J$ & gradient: vector of all $\partial J / \partial \theta_j$ \\
$\eta$ & learning rate (step size) \\
logit $z$ & the score before the sigmoid or softmax \\
\bottomrule
\end{tabular}
\caption{Learning and optimization notation summary.}
\label{tab:learning_notation}
\end{table}

\begin{table}[htbp]
\centering
\small
\begin{tabular}{ll}
\toprule
\textbf{Symbol} & \textbf{Description} \\
\midrule
$\sigma(z) = \frac{1}{1 + e^{-z}}$ & logistic sigmoid, values in $(0, 1)$ \\
$\zeta(a) = \log(1 + e^a)$ & softplus: a smooth version of $\max\{0, a\}$ \\
$e_y \in \{0, 1\}^k$ & one-hot vector: 1 in position $y$, 0 elsewhere \\
$\mathbb{I}[\text{cond}]$ & indicator: 1 if cond is true, else 0 \\
$\mathbf{1}$ & vector of ones \\
$\odot$ & element-wise product \\
\bottomrule
\end{tabular}
\caption{Functions and symbols notation summary.}
\label{tab:functions_notation}
\end{table}

\section{FeedForward Networks}
Also called multilayer perceptrons, a feedforward network is a composition of functions
with a certain \textbf{Depth}(number of layers with parameters), \textbf{Width}(number
of units $dI$ in layer $I$). The goal of training is to choose $\theta$ so that $$f(x; \theta) \approx f^*(x)$$
on new inputs, not only on training set. 
\subsection{Anatomy of a Unit and of a Layer}

\begin{tikzpicture}[
    input node/.style={circle, draw=green!40!black, fill=green!20, thick, minimum size=8mm, inner sep=0pt},
    neuron/.style={circle, draw=green!50!black, fill=green!60!black, text=white, thick, minimum size=1.2cm, inner sep=0pt},
    >=Stealth
]

    % Nodi di input
    \node[input node] (x1) at (0, 2) {$x_1$};
    \node[input node] (x2) at (0, 0) {$x_2$};
    \node[input node] (x3) at (0, -2) {$x_3$};

    % Neurone centrale
    \node[neuron] (neuron) at (3.5, 0) {$\sum \mid g$};
    \node[below=2pt of neuron, font=\small] {$+b$};

    % Archi con pesi
    \draw[thick, black!60] (x1) -- node[above left, font=\small] {$w_1$} (neuron);
    \draw[thick, black!60] (x2) -- node[above, font=\small] {$w_2$} (neuron);
    \draw[thick, black!60] (x3) -- node[below left, font=\small] {$w_3$} (neuron);

    % Uscita
    \draw[thick, black!60, ->] (neuron) -- ++(1.5, 0) node[right, black, font=\small] {$h$};

\end{tikzpicture}